\section{Budget-Aware Structured Control}

\subsection{Decision State}

At a decision point, the controller receives
\[
  z_t=(W_t,\mathbf{b}_t,\mathcal{H}_t),
\]
where \(W_t\) is the workspace, \(\mathbf{b}_t\) is the remaining resource
vector, and \(\mathcal{H}_t\) contains cost and outcome statistics available
from training or held-out traces. It first filters infeasible actions, including
any action that would consume the verifier call reserved for final assembly.
It then jointly ranks ready branches and applicable actions.

\subsection{Action Utility}

For state \(s\), branch \(n\), and action \(a\), we estimate:

\begin{itemize}
    \item \(\hat p_{\mathrm{succ}}\): probability that the action immediately
    yields an accepted artifact;
    \item \(\widehat{\Delta}\): expected verifier-grounded residual progress;
    \item \(U\): downstream unlock value, computed from obligations that would
    become ready if the target were accepted;
    \item \(I\): expected information gain over explicit competing failure
    hypotheses;
    \item \(R\): repetition and risk penalties, including unchanged goal
    fingerprints and mutation-scope risk;
    \item \(\hat c\): expected incremental scalarized cost.
\end{itemize}

The base score is
\[
  Q(s,n,a)=
  \frac{
    \alpha \hat p_{\mathrm{succ}}
    +\beta \widehat{\Delta}
    +\eta U
    +\kappa I
    -\rho R
  }{\epsilon+\hat c}.
\]
All terms are logged separately; the combined score is not treated as a
semantic truth. A heuristic controller can instantiate these terms from
diagnostic categories and graph structure, while a learned variant estimates
them from held-out traces. The main comparison must use no test-task outcomes
to fit either version.

\subsection{Conditioning on Remaining Budget}

Let \(r\in[0,1]\) be a normalized remaining-budget ratio after reserving final
assembly. We use a budget-conditioned concentration exponent
\[
  \gamma(r)=\gamma_{\max}
  -(\gamma_{\max}-\gamma_{\min})r,
\]
so abundant budget yields a flatter ranking and scarce budget concentrates
probability mass on high-value actions. The final priority is
\[
  \operatorname{Priority}(s,n,a)
  =\operatorname{sign}(Q)\lvert Q\rvert^{\gamma(r)}
  \cdot G(a,r),
\]
where \(G\) is an action-type gate. Early in a run, \(G\) permits
decomposition, capability tests, and representation changes that may create
future options. Near exhaustion it favors completion and repair on branches
with short estimated cost-to-proof, and suppresses structural actions whose
children cannot be completed within the residual budget.

For a vector budget there is no unique remaining fraction. We use the minimum
normalized slack over hard resources,
\[
  r=\min_j \frac{b_j-b^{\mathrm{reserve}}_j}{B_j},
\]
and reject actions that violate any component. We additionally report
alternative scalarizations as a robustness analysis.

\subsection{Diagnostic-Conditioned Actions}

Lean diagnostics change the conditional value of computation. Parser and local
type errors often favor repair; unknown identifiers favor retrieval or a
capability test; repeated unchanged goals lower the value of another identical
implementation attempt; and a sound-looking argument with persistent
implementation failures may favor a representation change. These rules
initialize the heuristic policy and define features for the learned policy.
They do not turn diagnostics into infallible mathematical diagnoses.

Competing failure hypotheses make information-seeking actions explicit. If a
cheap signature test can distinguish ``missing theorem'' from ``incorrect
argument,'' its value includes the avoided cost of pursuing the wrong route.
Infrastructure failures, by contrast, are handled by deterministic rules and
are not presented to the model as mathematical hypotheses.

\subsection{Two-Tier Model Escalation}

The controller treats model tiers as actions with different cost and success
models. A cheap model handles short implementation and repair attempts; a
strong model is available for detailed proofs, decomposition, or strategy
revision. Escalation is selected when its estimated success gain per cost
exceeds cheap alternatives, not after a fixed number of failures. To prevent
the learned score from starving the strong model until it is unusable, the
policy can reserve a configurable strong-call budget and expose an ablation
without that reserve.

\subsection{Controller Loop}

The structured loop is:

\begin{enumerate}
    \item initialize a workspace and seed all ready branches;
    \item reserve resources required for final assembly;
    \item enumerate type-valid actions on ready branches;
    \item filter actions that violate hard budgets or workspace invariants;
    \item score branch--action pairs using the remaining budget;
    \item execute the highest-priority computation and charge its realized
    cost;
    \item deterministically reduce the result into a successor workspace;
    \item if a complete compatible solution exists, assemble and recheck it;
    otherwise continue until no useful feasible action remains.
\end{enumerate}

On termination, the system returns an accepted proof or a structured partial
result containing accepted helpers, blocked/open obligations, preserved
alternatives, raw checker observations, and exact resource usage.
